\section{The Grammar of TFL}

So far, despite some basic symbolization as stand-ins for sentences, our
discussion has been entirely informal: our arguments are built out of
ordinary English sentences; our symbolizations are likewise parphrases of
English.

This week we will take our first steps to \emph{formal logic} by describing
the \emph{syntax} of Truth-Functional Logic (TFL). We will learn what we mean
by truth-functional when we discuss the \emph{semantics} of TFL. For the time-
being we will concentrate on the \emph{form} in formal.

First and foremost, we want to know what 'counts' as sentence in TFL. We can
start with a few definitions about the things that might be sentences:
\emph{expressions}.
\\
\begin{definition}[Basic TFL Expression]
  A \emph{basic expression} in TFL is any of the following symbols:
  \begin{enumerate}
  \item Atomic Sentence Letters: $A, B, C$
  \item Connectives: \enot, \eand, \eor, \eif, \eiff
  \item Brackets: (, ), [, ]
  \end{enumerate}
\end{definition}
Note that the atomic sentences can be made both \emph{enumerable} and
\emph{potentially infintite} with the addition of subscripts and/or
superscripts: $A_4, B^2, C_1^2$, etc.
\\
\begin{definition}[TFL Expression]
An \emph{expression} in TFL is any sequence of basic TFL expressions.
\end{definition}
Many of these expressions will be nonsense. So what we are really interested in
are the \emph{sentences} of TFL. There are potentially infinitely many of
these, so we will have to define the sentences \emph{recursively} using rules.
That is, we will need to explain how to construct sentences from expressions.
\\
\begin{definition}[TFL Sentence]\label{def:tfl-sentence}
  A \emph{sentence} in TFL is defined by the following rules:
  \begin{enumerate}
  \item Every atomic sentence letter is a sentence.
  \item If \ephi is a sentence, then \enot\ephi is a sentence.
  \item If \ephi and \epsi are sentences, then (\ephi\eand\epsi) is a sentence.
  \item If \ephi and \epsi are sentences, then (\ephi\eor\epsi) is a sentence.
  \item If \ephi and \epsi are sentences, then (\ephi\eif\epsi) is a sentence.
  \item If \ephi and \epsi are sentences, then (\ephi\eiff\epsi) is a sentence.
  \item Nothing else is a sentence.
  \end{enumerate}
\end{definition}
This is called an \emph{inductive} defininition: it uses rules to build infinitely many
complex expressions from a small number of basic one. Inductive defiitions form the
backbone of most systems of deductive logic, both in syntax and semantics. Notice that
the last clause ensures that there are many TFL expressions that are not TFL sentences.
But of course every TFL sentence is a TFL expression.

These rules allow us to identify a few important components of complex
sentences.
\\
\begin{definition}[Main Logical Operator]
  The \emph{main logical operator} of a sentence is the last connective applied in its
  construction.
\end{definition}
Given our definition of TFL sentences, we know that there is a unique main logical
operator for any (non-atomic) sentence. Looked at from the other direction, we can
say that the \emph{scope} of a connective is just the new sentence that its rule makes.
\\
\begin{definition}[Connective Scope]
  The \emph{scope} of a connective is the sentence for which it is the main logical
  operator.
\end{definition}
When checking to see whether some expression is a sentence, we will need to work
backward through these rules: break the expression apart into component parts until we
are left with only atomic sentence letters. Doing this is usually straight-forward and
involves finding the main logical operator (if any) in the expression under review.
Indeed, our bracketing conventions make the main logical operator and its scope almost
immediately apparent. But as we get better at reading complex expressions, we will treat
these brackets more and more as optional or only occasionally obligatory. For the time
being, we will adopt the conventions in \cite[][]{Mag25} by omitting outermost brackets
and sometimes substituting [, and ] for clarity and readability.

In ordinary English, we can find two distinct types of \emph{ambiguity}. \emph{Lexical
ambiguity} is when on word can mean more than one thing. A popular example is the word
\emph{bank} which can be either a financial institution or the edge of a river. \emph{
Structural ambiguity}, on the other hand, is when a sentence can have more than one
meaning based on how the parts work together. A popular example is the sentence
`The astronomer saw his student with a telescope', which can have two different readings:
one where the astronomer has the telescope and one where the student does. A closely
related concept is \emph{vagueness} which describes terms that are imprecise and often
difficult to apply.

As intimated above, TFL sentences are \textbf{never} structurally ambiguous. And we can
make efforts in formulating our symbolization keys to minimize, if not elminate, both
vagueness and lexical ambiguity. That said, all three forms of ambiguity are rampant in
ordinary language and our efforts at symbolizing complex arguments from English will
need to proceed with due care. Unlike TFL, English does not have bracketing rules. We
will refer to the structural ambiguities that can give more than one TFL sentence as
\emph{scope ambiguities}. Given the sentence `The professor is literate or not
dumb and lazy' and the symbolization key:

\begin{argument}
\item[A:] The professor is literate.
\item[B:] The professor is dumb.
\item[C:] The professor is lazy.
\end{argument}

We can identify a few scope ambiguities, offering plenty of opportunities for the
professor to be illiterate, dumb or lazy but not, as it turns out, all three. Notice
how the bracketing helps to distinguish between the different \emph{readings} of the
English sentence being symbolized.

\begin{argument}
  \item[] $(A \eor \enot B) \eand C$
  \item[] $A \eor \enot (B \eand C)$
  \item[] $A \eor (\enot B \eand C)$
\end{argument}

In our informal discussion, it was convenient for us to mix together the sentences of
the logic we were discussing and the sentences about that logic. Now that we are dealing
with formal logic, we will need to be a bit more careful. We will also need to be clear
about when we are using English in an ordinary way versus when we are using it to talk
about a language, whether English or TFL. To this end, a few more conventions and
definitions are in order.

The most common indicator of level-mixing is quotation marks. Consider the following
examples.

\begin{argument}
\item[A:] `Aristotle' is not the original Greek spelling.  
\item[B:] Aristotle taught Alexander the Great.
\end{argument}

Whereas B is about an English word, A is about the ancient philosopher. This is an
importnant distinction: the former is called \emph{use}; and the latter is called
\emph{mention}. Speaking about the phenomenon more generally, we can adopt the following
definitions.

\begin{definition}[Object Language and Metalanguage]
  When talking about TFL (or English), we call the language that we are talking
  about the \emph{object language} and the language that we are using to talk
  about it the \emph{metalanguage}.
\end{definition}

Since TFL and our metalanguage do not use quotes, we will rarely need to worry about
them. The most important point to make here is that we will use English as our
metalanguage when discussing TFL (and later FOL). But we will make some adjustments to
simplify our presentation. As we have already seen in Definition-\ref{def:tfl-sentence},
we will use Greek lower case letters `\ephi', `\epsi', `\echi', etc as
\emph{metavariables} for arbitrary TFL sentences. Upper case Latin letters $`A', `B',
`C'$, etc will continue to be used as atomic sentence letters. After now, we will rarely
mention these variables, even in the metalanguage. So the quotation marks will almost
always be omitted.

One more special feature of our English++ metalanguage is the conclusion marker
`\therefore', which we've seen already in the informal discussion of arguments. We can
use someshorthand to formalize the notation in our metalanguage as well. For instance,
we can describe a sequence of premises and conclusion that comprise an argument as:

\begin{argument}
\item[] $\phi_1,\dots,\phi_n  \therefore  \psi$
\end{argument}

But we will try to strike a healthy balance in the metalanguage between too much verbiage
and too much formalism.

